\section{Related Work}
\label{sec:related}

\subsection{Classification from reduced and missing leads}
The question of how much of the twelve-lead ECG a classifier really needs was put to the community by the
PhysioNet/Computing in Cardiology Challenge 2021, \emph{Will Two Do?}, which asked whether reduced-lead
recordings retain the diagnostic information of the full twelve leads \cite{reyna2021willtwodo}. The winning
entry, an ensemble of residual convolutional networks with a multi-head attention mechanism, ranked first in every
lead configuration, with nearly unchanged scores from twelve down to two leads on the hidden validation set
\cite{nejedly2021cinc,nejedly2022ensemble}. In representation learning, Oh \emph{et al.}\ pre-train ECG encoders with random lead
masking so that the learned features do not depend on a fixed set of leads \cite{oh2022leadagnostic}.

These works share a premise: the model is told, or trained to expect, that some leads are absent. Our setting
differs in two ways. First, failure is unannounced: the input is always a twelve-lead array, and a failed lead
arrives as zeros, a flat line or noise. Second, our aim is not a better reduced-lead classifier but a controlled
attribution: given that training already includes lead failure, as in lead masking, does an explicit per-lead
quality estimate that gates the attention add anything? Our answer, that the augmentation carries almost all of
the benefit, is consistent with the success of masking-based training in this prior work.

\subsection{Benchmarks on PTB-XL}
Strodthoff \emph{et al.}\ established the reference benchmark for PTB-XL with its recommended patient-level folds;
convolutional networks, in particular ResNet- and Inception-based architectures, performed best across tasks
\cite{strodthoff2021benchmark}. We adopt the same folds, label definition and 100\,Hz signal. Our models are
trained under a deliberately small, shared budget and do not reach the benchmark's absolute scores
(Section~\ref{sec:res-clean}); every comparison in this paper is therefore relative, between models trained
identically.

\subsection{Evidential deep learning and its critics}
Evidential deep learning places a Dirichlet distribution over class probabilities, interprets the network
outputs as evidence in the sense of subjective logic, and reads uncertainty from the total amount of evidence,
all from a single deterministic network \cite{sensoy2018edl}. Its original evaluation reported strong
out-of-distribution detection and robustness to adversarial perturbation. Its theoretical footing has since been
questioned. Bengs \emph{et al.}\ showed that the loss functions used to train such second-order predictors do not
incentivise them to represent their epistemic uncertainty faithfully \cite{bengs2022pitfalls}. Shen
\emph{et al.}\ sharpened this analysis across many evidential objectives, found that the learned epistemic
uncertainty does not vanish even with infinite data, and argued that evidential models behave more like
energy-based out-of-distribution detectors, which can still perform well on downstream tasks despite poor
uncertainty quantification \cite{shen2024mirage}.

Our results add a concrete, applied data point to this debate. Lead failure is a structured loss of information
whose severity we control exactly, and the evidential uncertainty of our model does not track it
(Section~\ref{sec:res-edl}). Unlike in an out-of-distribution benchmark, a recording with failed leads
contains strictly less information than the same recording intact, so an uncertainty measure that stays flat
as leads disappear is hard to defend.

\subsection{Calibration and selective prediction}
Modern networks are often poorly calibrated, and a single temperature fitted on held-out data is a simple and
effective remedy \cite{guo2017calibration}; we apply it to every sigmoid model so that calibration comparisons
with the evidential head are fair. Selective classification lets a classifier abstain to reach a target risk;
thresholding the network's own maximum class probability is already a strong rejection rule
\cite{geifman2017selective}. Our confidence score is the multi-label analogue of that rule, and our
risk--coverage results show that it matches the evidential uncertainty (Section~\ref{sec:res-edl}).

\subsection{What this study adds}
We do not propose a new architecture. We contribute a controlled ablation that separates three ingredients that
are usually introduced together (training with lead failure, quality-gated lead attention and evidential
abstention), evaluates them on identical, deterministically corrupted test data, controls for the accuracy of
the quality estimate with an oracle, and reports paired confidence intervals. The value of the study lies in
which of these ingredients survive that separation.
